% source: https://github.com/pfradin/luadraw/discussions/140#discussioncomment-15060325 (pfradin, block 1)
\documentclass[boder=5pt]{standalone}
\usepackage[svgnames]{xcolor}
\usepackage[3d]{luadraw}
\usepackage{fourier-otf}
\begin{document}
\begin{luadraw}{name=tubes}
local g = graph3d:new{window={-6,6,-4,7},bg="lightgray!30" }
local O, r = Origin, 0.5
local R  = r*math.sqrt(2)
local S = sphere(O,R)
local C = parallelep(M(-r,-r,-r), 2*r*vecI,2*r*vecJ,2*r*vecK)
local crossing_base = clip3d(S,C)
local Xtube = poly2facet( cylinder(-r*vecI,r,r*vecI,20,true) )
local Ytube = poly2facet( cylinder(-r*vecJ,r,r*vecJ,20,true) )
local Ztube = poly2facet( cylinder(-r*vecK,r,r*vecK,20,true) )

local crossing = function(center)
    return shift3d(crossing_base,center)
end

local tube = function(center, dir) -- dir = "x" or "y" or "z"
    if dir == "x" then return shift3d(Xtube,center)
    elseif dir == "y" then return shift3d(Ytube,center)
    elseif dir == "z" then return shift3d(Ztube,center)
    end
end

local grid_base = {}
for y = -4, 4  do
    local A = y*2*r*vecJ
    if y%3 == 0 then
        insert(grid_base, crossing(A))
        insert(grid_base, tube(A+2*r*vecI,"x"))
        insert(grid_base, tube(A-2*r*vecI,"x"))
        insert(grid_base, tube(A+2*r*vecK,"z"))
        insert(grid_base, tube(A-2*r*vecK,"z"))
    else
        insert(grid_base, tube(A, "y"))
    end
end
local k = 6*r
local grid = concat(grid_base,shift3d(grid_base,k*vecI),shift3d(grid_base,-k*vecI))
grid = concat(grid, shift3d(grid,k*vecK))
g:Dfacet(grid ,{color="orange",mode=mShadedOnly, contrast=0.5})
g:Show()
\end{luadraw}
\end{document} 
